% Mathematical proof component for AX11H and QAX11H.
% Adapted with attribution from the delivered AX11 preservation argument and
% PA194Q's exchange argument; the cached-only AX11 Pop/MP contracts and the
% bounded-input induction below are the changes checked in this review.
% Requires amsmath, amssymb, amsthm, booktabs.
\section{Preservation for the eleven-entry cached-test composition}
\label{sec:ax11h-preservation}

This section composes the eleven-entry, four-parameter AX interface of
the delivered PA194AX11 packet with the cached-only Shift test. An optional
third component exchanges the two premises in the quantifier-distribution
head. The command equations below are \emph{not} those of the full PA194Q
machine: that machine also changes successful MP to return its result in
the cache. Here successful MP returns its result in the visible top and
does not push an extra stack entry.

We call the two abstract command systems $R_{11H}$ and $R_{Q11H}$.
The first searches the inherited separate-diagonal numerical calculus
$\mathcal C_4$. The second searches its premise-exchanged version
$\mathcal C_4^Q$. Any assertion about a delivered physical table also
requires its separate exact source, runtime, and PA calibration proofs.
The subsequent sections supply those interfaces and compose the full
theorem in Section~\ref{sec:complete-equivalence}.

\subsection{The two numerical calculi}

Use Cantor pairing and the inherited coefficient-four constructors
\[
 p(a,b)=\frac{(a+b)(a+b+1)}2+b,\qquad
 I(A,B)=4p(A,B),\quad M(a,b)=4p(a,b)+1,
\]
\[
 U(a,B)=4p(B,a)+2,\qquad E(a,b)=4p(a,b)+3.
\]
Write $A\to B$ for $I(A,B)$, $N(B)=I(B,1)$,
$A\wedge B=N(I(A,N(B)))$, and
$\exists_a B=N(U(a,N(B)))$. Zero is the numerical truth code;
one is the designated contradiction. In particular $I(0,0)=0$.

The unchanged ordinary heads are
\begin{align*}
 h_1(D,H)&=(\exists_D U(D,H))\to H,\\
 h_2(D,A,H)&=U(D,U(A,H))\to U(A,U(D,H)),\\
 h_3(D,A,H)&=M(H,D)\to(E(H,A)\to M(A,D)),\\
 h_4(D,A,H)&=E(H,D)\to(M(A,H)\to M(A,D)),\\
 h_5(D,H)&=\exists_D E(H,D).
\end{align*}
The two choices for the sixth head are
\begin{align*}
 h_6(D,A,H)&=U(D,A)\to(U(D,I(A,H))\to U(D,H)),\\
 h_6^Q(D,A,H)&=U(D,I(A,H))\to(U(D,A)\to U(D,H)).
\end{align*}
All other axiom families are identical:
\begin{align*}
 \Gamma(B)&=I(B,U(B,B)),\\
 \rho_8(D,A,H,u)&=((H\to u)\to D)\to((D\to H)\to(A\to H)),
 \quad u\leq D,
\end{align*}
and the guarded gate
\begin{align*}
 V(a,b,c)&=M(a,b)\to U(c,M(a,b)),\\
 P(a,b,c)&=c\to M(c,b),\\
 S(a,b,c)&=U\bigl(b,(M(a,b)\to P(a,b,c))\wedge(M(c,b)\to c)\bigr),\\
 Q(a,b,c)&=U(c,S(a,b,c)\to N(P(a,b,c))),\\
 G(a,b,c)&=N(V(a,b,c)\to Q(a,b,c)),\qquad c\ne a,b.
\end{align*}
The gate emitter returns zero if its guard fails. Modus ponens and
theorem generalization $B\mapsto U(i,B)$ are the two inference rules.
The fixed leaf $\rho_8(0,0,0,0)=0$ supplies a proof of zero.

The inherited restricted implicational basis derives premise exchange
\[
 (a\to(b\to c))\to(b\to(a\to c))
\]
without either choice of the sixth head. Consequently the two sixth-head
families are interderivable uniformly in $D,A,H$. Replacing such axiom
leaves and retaining MP and generalization gives the inherited
primitive recursive proof translations between
$\mathcal C_4$ and $\mathcal C_4^Q$, justified in iPRA in
Section~\ref{sec:q-exchange-ipra}. This is exactly the logical exchange
argument of the PA194Q packet; the present composition does not count it
as a new mathematical idea. It also does not import that packet's cached
MP command semantics.

\subsection{The exact four-parameter AX effect}

Let $[a;w]=p(a,w)$, with the first displayed entry at the stack top.
Cantor projections are total:
\[
 \operatorname{hd}[a;w]=a,\quad
 \operatorname{tl}[a;w]=w,\quad
 \operatorname{hd}(0)=\operatorname{tl}(0)=0.
\]
Thus $[0;0]=0$ is an allowed numerical alias. None of the following
calculations treats these stacks as free lists.

The store is $(t,w,d,h)$. Set $D=d+x$, $H=h+y$, and $A=c$. The
eleven-entry command has the effect
\[
 \mathsf A_{c,x,y,z}(t,w,d,h)
   =(0,[q_0,q_1,\ldots,q_9,t;w],z,0),                  \tag{H-A}
\]
where
\[
\begin{array}{c|l}
 j&q_j\\\hline
 0&G(H,D,A),\text{ or zero when the guard fails}\\
 1&\Gamma(A)\\
 2&h_6(D,A,H)\text{ or }h_6^Q(D,A,H)\\
 3&h_5(D,H)\\
 4&h_4(D,A,H)\\
 5&h_3(D,A,H)\\
 6&h_2(D,A,H)\\
 7&h_1(D,H)\\
 8&\rho_8(D,A,H,d)\\
 9&U(d,h).
\end{array}
\]
The extra operand $z$ is real. The AX tail pushes its selected gate and
falls into the existing full MP field reader; that reader consumes
another continuation field. Immediately after the push the visible
values are $t=d=h=0$. The ensuing MP compares $I(0,z)$ with zero.
For $z=0$ it succeeds, and for $z>0$ it fails; both branches leave
$t=h=0$, $d=z$, and the newly extended stack unchanged. Therefore (H-A)
holds for every $z$, not merely for the $z=0$ calls used by the compiler.

All $q_0,\ldots,q_8$ have residue zero modulo four.
$q_9=U(d,h)$ has residue two and is at least two. In particular no
generated entry is one. The exchange in $q_2$ preserves this fact and
all stack positions.

The other two commands are exactly
\[
 \mathsf S(t,w,d,h)=
 \begin{cases}
  \mathsf{halt},&h=1,\\
  (h,\operatorname{tl}w,d,\operatorname{hd}w),&h\ne1,
 \end{cases}                                             \tag{H-S}
\]
and, for $b=d+x$,
\[
 \mathsf{MP}_x(t,w,d,h)=
 \begin{cases}
  (b,w,b,0),&I(t,b)=h,\\
  (0,w,b,0),&I(t,b)\ne h.
 \end{cases}                                             \tag{H-M}
\]
Only $\mathsf S$ halts. In particular $t=1,h=0$ does not halt under
this Shift, whereas $t=h=1$ does. These are the opposite outcomes of
the AX11 sum guard in those two cases. A preservation proof must use
(H-S), not silently reuse the old $t+h=1$ test.

\subsection{Local word equations}

Commands execute from left to right. Put $A_0=\mathsf A_{0,0,0,0}$.
Successive applications of (H-S) and the projection equations give
\begin{align*}
 \mathsf A_{c,x,y,0};\mathsf S^{10}(t,w,d,h)
    &=(q_8,[t;w],0,q_9),\\
 \mathsf A_{c,x,y,0};\mathsf S^{11}(t,w,d,h)
    &=(q_9,w,0,t).
\end{align*}
Both executions are total: their tests see the cleared cache zero and
then $q_0,\ldots,q_9$, never the saved top. Define
\[
 \mathsf{Push}=A_0;\mathsf S^{10};\mathsf{MP}_0,\qquad
 \mathsf{Pop}=A_0;\mathsf S^{12},\qquad
 \mathsf{Discard}=\mathsf{MP}_0,\qquad
 \mathsf{Offset}_i=\mathsf{MP}_i.
\]
The exact contracts are
\begin{align*}
 \mathsf{Push}(t,w,d,h)&=(0,[t;w],0,0),\\
 \mathsf{Pop}(t,[a;w],d,h)
   &=\begin{cases}
       \mathsf{halt},&t=1,\\
       (t,w,0,a),&t\ne1,
      \end{cases}\\
 \mathsf{Discard}(t,w,0,h)&=(0,w,0,0),\\
 \mathsf{Offset}_i(0,w,0,0)&=(0,w,i,0).
\end{align*}

Push's final comparison is between an implication code, residue zero,
and the universal code $q_9$, residue two, so it fails. For Pop, the
twelfth Shift sees the saved top in its cache and therefore halts
\emph{exactly} when that top is one. This is the essential change from
the original AX11 Pop, which was unconditional. Discard's two branches
coincide when its proposed conclusion is zero. Offset at $i=0$ has
coinciding zero outputs; at $i>0$, $I(0,i)>0$ forces its failure branch.
These calculations include all zero aliases.

For $0\le j\le8$, the selector
\[
 \mathsf{Discard};\mathsf A_{A,D,H,0};
                   \mathsf S^{j+2};\mathsf{Pop}^{9-j}       \tag{H-Sel}
\]
sends every $(t,w,0,0)$ to $(q_j,w,0,0)$. After the Shifts, its
top is $q_j$, its cache is $q_{j+1}$, and the remaining entries above
the base are $q_{j+2},\ldots,q_9,0$. There are exactly $9-j$ such
entries. Each Pop preserves the same $q_j\ne1$, removes one entry,
and places that entry in the cache. The last removed entry is the
saved zero, so the final cache is zero. In particular the gate
selector uses $j=0$, the diagonal selector $j=1$, and the sixth-head
selector $j=2$, regardless of the premise exchange.

For the restricted leaf $\rho_8(D,A,H,u)$, $u\le D$, use
\[
 \mathsf{Discard};\mathsf{Offset}_u;
 \mathsf A_{A,D\mathbin{\dot{-}}u,H,0};
                      \mathsf S^{10};\mathsf{Pop}.          \tag{H-Rho}
\]
The old offset is $u$, the accumulated parameter is $D$, and the
fourth operand zero clears the offset before the selector finishes.
The preserved restricted head has residue zero, so its Pop is safe.

Generalization is
\[
 \mathsf{Gen}_i=\mathsf{Push};\mathsf{Offset}_i;\mathsf S;
                 \mathsf A_{i,0,0,0};\mathsf S^{11}.         \tag{H-Gen}
\]
It sends $(a,w,0,0)$ to $(U(i,a),w,0,0)$ for every $a,w,i$,
including $a=1$. Before the middle Shift the cache is zero; that
Shift moves $a$ into the cache. The following AX consumes it as its
generalization parameter and clears it before another Shift tests it.
The eleven-Shift identity then returns $U(i,a)$ with zero cache.

All local identities are assertions about primitive recursive command
functions and follow in standard iPRA from their defining equations,
Cantor projection equations, the displayed residue facts, and finite
case splits on decidable numerical equalities. This statement concerns
the abstract commands, not the literal machine's implementation of them.

\subsection{Modus ponens and an all-store compiler}

From $(0,[a,I(a,b),0;w],0,0)$ the suffix
\[
              \mathsf S;\mathsf S;\mathsf{MP}_b;\mathsf{Pop}       \tag{H-MPs}
\]
halts when $a=1$ or $b=1$, and otherwise returns $(b,w,0,0)$.
The first Shift retrieves $a$ into the cache. The second either
halts at $a=1$ or makes $a$ the top and $I(a,b)$ the cache.
MP succeeds and returns $(b,[0;w],b,0)$. Pop clears the offset,
removes the reserved zero, and either halts at $b=1$ or returns
$(b,w,0,0)$. The reserved zero is essential; without it an
underlying stack entry would be consumed. The calculation remains
valid at $a=b=w=0$ because it uses total projection identities.

Compile numerical proofs into words. Axiom leaves use (H-Sel) or
(H-Rho); a generalization appends (H-Gen). At an MP node with major
proof $e_1$ of $I(a,b)$ and minor proof $e_0$ of $a$, set
\[
 W_e=\mathsf{Discard};\mathsf{Push};W_{e_1};\mathsf{Push};
        W_{e_0};\mathsf{Push};\mathsf S;\mathsf S;
                                  \mathsf{MP}_b;\mathsf{Pop}.
\]
The exact invariant is
\[
 W_e(t,w,0,0)=\mathsf{halt}
 \quad\hbox{or}\quad
 W_e(t,w,0,0)=(\operatorname{concl}(e),w,0,0).             \tag{H-Cmp}
\]
The first Discard and Push in the MP case reserve the zero. The
major computation and its Push put the major premise over that
zero; the minor computation and its Push put the minor premise
over both. The suffix's displayed precondition follows exactly.
An earlier halt satisfies the invariant and is preserved by
every word extension.

For a refutation $e$, the word
\[
                        A_0;W_e;\mathsf{Pop}               \tag{H-Ref}
\]
halts from every incoming store $(t,w,d,h)$. Its first AX normalizes
$d=h=0$ and produces some top and finite stack; (H-Cmp) applies
without any empty-store assumption. If the compiled proof has not
already halted, it returns top one with zero cache; Pop then tests
that saved one and halts. A terminal plain Shift would fail: it
would see cache zero. Appending a Shift after Pop is harmless but
unnecessary, since the refutation already halts within Pop.

\subsection{A quantifier-free proof of the compiler invariant}

Ordinary induction on a statement containing ``for every base
stack'' is not automatically an allowed induction in iPRA. Here is
an explicit bounded reformulation for this particular compiler.
It proves (H-Cmp) in standard iPRA for the abstract command functions.

Encode a proof as a finite sequence of $N>0$ lines; every MP or
generalization premise is an earlier line. The primitive recursive
proof checker tests those index inequalities, the numerical axiom
families, and the required conclusion equalities. Compile $W_j$
bottom-up for each line $j<N$. This is a primitive recursion on
$j$: look up the previously stored words and concatenate the finite
prefixes and suffixes displayed above. On malformed data use an
arbitrary fixed word. No normalization or proof search occurs in
this compilation.

Let $L$ be the maximum length of these $N$ compiled words and $K$
the maximum command code appearing in them. Both are primitive
recursive functions of the proof code. Choose a total primitive
recursive command step $\operatorname{next}(c,u)$ on coded stores,
with an absorbing halt flag, and put
\[
 \Phi_K(s)=1+\max\left(s,
    \max_{c\le K,\ u\le s}\operatorname{next}(c,u)\right).
\]
This is a primitive recursive, monotone upper bound for one command
step. Its bounded maxima are computed by primitive recursion; there
is no unbounded maximization. Every prefix of length at most $L$ of
any $W_j$, started at a store code at most $s$, stays below
$\Phi_K^L(s)$. Quantifier-free induction proves this elementary
iteration bound.

Fix an arbitrary normalized initial store code $s_0$ and define
\[
 B(j)=\Phi_K^{\,L(N-1-j)}(s_0),\qquad j<N.
\]
If $i<j$, then
\[
                  \Phi_K^L(B(j))\le B(i).                 \tag{H-Bnd}
\]
Prove successively for every $j<N$ that (H-Cmp) holds for every
normalized input store of code at most $B(j)$. This bounded
universal predicate has a primitive recursive characteristic
function: evaluate the finite word, check whether the result is
halted or the specified store, and take a bounded conjunction over
all store codes up to $B(j)$. More explicitly, the induction invariant
at stage $j$ is the bounded conjunction of these correctness tests for
every $i<j$ and every normalized store of code at most $B(i)$.
It retains every earlier premise, not only line $j-1$. The induction
over this finite prefix is quantifier-free after standard Boolean
coding.

For an axiom line use the local selector identity. For
generalization or MP, each premise word is called after a prefix
of at most $L$ commands of the current word. Unless execution
has already halted, its actual input is normalized by the preceding
local contracts and has code at most $\Phi_K^L(B(j))$. By (H-Bnd)
it lies within the induction hypothesis for its earlier premise
line. The local equations now establish the conclusion. This
completes the quantifier-free induction.

At the last line $B(N-1)=s_0$, so the theorem applies to the
arbitrary original store. Universal generalization yields
(H-Cmp) for every proof and every normalized store; prefixing AX
gives (H-Ref) for every store. The bounds are deliberately huge.
They establish primitive recursiveness and allowed induction
strength, not an efficient execution algorithm. This argument
does not establish the unrelated PA-to-numerical-calculus compiler.

\subsection{Four-field word encoding and exhaustive scheduling}

Let $e(\epsilon)=0$ and $e(c::w)=1+b_c(e(w))$, with
\begin{align*}
 b_{\mathsf S}(z)&=p(0,z),\\
 b_{\mathsf{MP}_x}(z)&=p(1,p(x,z)),\\
 b_{\mathsf A_{c,x,y,u}}(z)&=p(c+2,p(y,p(x,p(u,z)))).       \tag{H-Enc}
\end{align*}
The field order is $y,x,u$ after the AX tag. Total projections
give the inverse equations even at zero. Retaining a three-field
AX encoder would make the new MP reader consume a piece of the
following continuation; it is not an equivalent encoding.

An exact live scheduler configuration is $(n,\ell,t,w,d,h)$.
The clean initial configuration is $(2,0,0,0,0,0)$. Its next step
selects a raw seed and updates the next-unused fresh counter by
\[
 (\text{seed},n')=
 \begin{cases}
 (n,n+1),&\ell=0,\\
 (\ell-1,n),&\ell>0.
 \end{cases}
\]
It decodes and executes one command, sets the new continuation to the
residual after all that command's fields, and retains $n'$. Thus the
fresh counter advances when a candidate is selected. The store is
carried between candidates. A halted configuration is absorbing.

The continuation produced by any of the
three decoders is a nested Cantor component of its raw seed. Thus
it is at most that seed, and for a positive continuation it
strictly decreases: the successor continuation is at most
$\ell-1$. The extra AX field preserves this elementary bound.

Under this fetch equation, a fresh candidate $n$ has at most
$n+1$ command invocations before it either halts or returns to
the fresh-candidate boundary. Every abstract command is a total
primitive recursive operation, with a halt flag. Hence a
specified candidate $N\ge2$ is reached from the clean initial
candidate two within a primitive recursive finite bound, unless
the scheduler halts earlier. For the command-step convention just
described, $(N+1)^2$ is a sufficient loose bound for draining the
earlier candidates; a different administrative-step convention
only adds its explicitly bounded overhead.

Given the all-store halting suffix (H-Ref), prefix another AX
command $\mathsf A_{k,0,0,0}$. Its fresh raw seed is
\[
 p\bigl(k+2,p(0,p(0,p(0,e(\text{suffix}))))\bigr).        \tag{H-Cof}
\]
These seeds are unbounded in $k$ and at least three. The prefix
may alter the carried store; the suffix halts from every store.
Since the displayed seed is at least $k+2$, taking $k$ to be the
current next-unused counter is an explicit cofinal choice. It
provides a finite halting witness. In
particular, from the clean candidate-two start one may choose
$k=0$. There is no assumption that failed trials reset the stack.

The decoding identities, strict-decrease bound, and finite
candidate bound are primitive-recursive arithmetic statements
proved by quantifier-free induction. This yields an iPRA
completeness theorem for the explicit abstract scheduler with
these fetch equations and the commands (H-A),(H-S),(H-M).
Identifying a literal source program with those equations is a
separate binding obligation, just as identifying it with a
physical binary machine is separate.

\subsection{Soundness by finite proof annotations}

Use the appropriate calculus $\mathcal C_4$ or $\mathcal C_4^Q$.
Maintain a finite valid derivation list together with indices
proving the current top $t$, cache $h$, and every canonical
Cantor-stack head. The fixed zero proof supplies virtual zero
entries when the numerical stack is zero. For $w>0$,
$\operatorname{tl}(w)<w$, so canonical stack decoding can be
bounded by $w$ iterations; the associated annotation invariant
has a primitive recursive Boolean checker.

AX appends proofs of its generated entries. Its saved top already
has a proof; $q_9=U(d,h)$ follows by generalization from the cache;
$q_8$ satisfies the actual guard $d\le d+x$; the other entries
are its selected calculus axioms or the fixed zero in a failed
gate case. The fallthrough MP puts zero in the two visible
positions for every fourth operand. The offset $d$ carries no
derivability requirement.

Successful MP combines the proofs of $t$ and $I(t,b)=h$ to obtain
$b$; failure returns the fixed zero proof. A nonhalting Shift
moves the cache proof to the top and the first stack annotation
to the cache, or uses the zero proof for an empty numerical stack.
A halting Shift has $h=1$ and therefore exposes a finite
refutation. The top need not be inspected for this soundness
direction.

Each annotation update is primitive recursive on the store and
the finite derivation list. The local validity checks and the
canonical zero-alias cases prove that its Boolean invariant is
preserved. Quantifier-free induction on the number of abstract
commands then extracts a calculus refutation from every finite
scheduler halt. This is a syntactic proof extraction, not a
truth interpretation of arbitrary formulas.

Together with completeness, standard iPRA proves for the explicit
abstract schedulers
\[
 \operatorname{Halt}(R_{11H})\leftrightarrow
          \exists e\,\operatorname{Prf}_{\mathcal C_4}(e,1),
\]
\[
 \operatorname{Halt}(R_{Q11H})\leftrightarrow
          \exists e\,\operatorname{Prf}_{\mathcal C_4^Q}(e,1),
\]
with the primitive recursive encodings and exact abstract step
equations just specified. These are written mathematical
weak-arithmetic arguments; no new Lean or Rocq build is asserted.

\subsection{Composition with the execution and PA proofs}

The remaining sections supply the two execution interfaces used by this
abstract proof. Section~\ref{sec:physical-ipra} proves the literal binary
table's correspondence with its exact primitive-register source, including
uniform rule soundness, frame closure, and complete initialization.
Section~\ref{sec:source-command-complete} matches every primitive source
instruction to explicit arithmetic templates and composes their exact
primitive recursive clocks through fetch, dispatch, and command execution.
The positive-segment transport theorem in
Section~\ref{sec:ipra-transport} therefore identifies each literal machine's
halting with refutability in its corresponding numerical calculus.

Section~\ref{sec:consistency-nonhalting} combines those proofs with the
pinned model-uniform numerical seed soundness and the one-way
$\mathrm I\Sigma_1$ transfer. It proves
$\mathrm{iPRA}\vdash\Con(\PA)\to\neg\Halt_{\rm blank}(M)$ for both
new literal tables. The converse uses the explicit Foundation adapter
and finite-demand logical compiler in
Sections~\ref{sec:foundation-paeq-complete}--\ref{sec:forward-calibration-completion}.
Section~\ref{sec:complete-equivalence} combines the two directions.
The bounded command-word argument above is one component of that
composition; the logical compiler has its own complete construction.
